\subsection{Zeitdarstellung in natürlicher Sprache}
Die Darstellung von Zeit erfolgt in technischen Systemen üblicherweise in numerischer Form, beispielsweise im Format \textit{HH:MM}.
Diese Form basiert auf einer direkten Abbildung von Stunden- und Minutenwerten.
Demgegenüber steht die sprachliche Zeitdarstellung, bei der Zeitangaben durch natürliche Sprache ausgedrückt werden.
Diese Form orientiert sich an linguistischen Strukturen und alltagssprachlichen Konventionen~\cite[S.~14--35]{klein1994time}.

Die sprachliche Darstellung der Uhrzeit im Deutschen folgt keiner strikt linearen Struktur.
Zeitangaben werden häufig relativ zu einem Referenzpunkt formuliert.
Insbesondere erfolgt eine Orientierung an der nächsten vollen Stunde.
Ein typisches Beispiel ist die Angabe „halb drei“, die einer Zeit von 2:30 entspricht und somit auf die folgende Stunde referenziert.
Diese Form der Darstellung unterscheidet sich grundlegend von der numerischen Repräsentation, bei der die aktuelle Stunde als Bezugsgröße dient~\cite[S.~36]{klein1994time}~\cite[S.~xv]{lakoff1987women}.

Ein weiteres Merkmal ist die Verwendung diskreter Zeitintervalle.
Die sprachliche Zeitdarstellung erfolgt typischerweise in Schritten von fünf Minuten, beispielsweise „fünf nach“, „viertel vor“ oder „zwanzig nach“.
Zur genaueren Beschreibung können zusätzliche Feinauflösungen genutzt werden, etwa durch Formulierungen wie „kurz nach“ oder „fast“.
Diese Struktur stellt eine abstrahierte und zugleich alltagstaugliche Repräsentation von Zeit dar~\cite[S.~208\,ff.]{langacker1987foundations}.

Darüber hinaus existieren regionale Unterschiede in der sprachlichen Zeitdarstellung.
Im deutschsprachigen Raum variieren insbesondere die Ausdrücke für Viertelstunden.
Während in einigen Regionen die Formulierung „Viertel nach zwei“ verwendet wird, ist in anderen Regionen „Viertel drei“ gebräuchlich, um denselben Zeitpunkt zu beschreiben.
Solche Unterschiede erfordern bei der technischen Umsetzung eine eindeutige Festlegung der zugrunde liegenden Sprachlogik~\cite[S.~212]{kucharczyk2012sprachliche}.

Aus kognitionswissenschaftlicher Perspektive bietet die sprachliche Zeitdarstellung Vorteile hinsichtlich Verständlichkeit und intuitiver Verarbeitung.
Alltagssprachliche Strukturen entsprechen erlernten kognitiven Mustern und ermöglichen eine schnelle Interpretation ohne explizite Umrechnung numerischer Werte.
Sprachbasierte Repräsentationen sind eng mit menschlichen Denkprozessen verknüpft und unterstützen eine effiziente Informationsverarbeitung~\cite[S.~15--18]{pinker1994language}~\cite[S.~xv]{lakoff1987women}.